\begin{table}[ht]
\centering
\caption{Samples and identifying variation.}
\label{tab:sample_variation_v9}
\begin{threeparttable}
\small
\setlength{\tabcolsep}{5pt}
\begin{tabular}{@{}l l l l@{}}
\toprule
Sample & Size & Unit & Role \\
\midrule
Full BEC pharmaceutical file & \BPnPOIfull{} & POI & Linked POI analysis file \\
Analysis sample & \BPnAnalysisSample{} & POI & Ordinary-vs-urgent comparisons \\
Winners-only price sample & \BPnegNobs{} & Winning bids & Negotiated-price regressions \\
Urgent panel & \BPnUTG{} & Winning bids & Admin-vs-litigated comparison \\
Both-regime urgent cells & \BPbothCellN{} & Item-month cells & Within-cell variation \\
Singleton urgent cells & \BPbothCellNSingleton{} & Item-month cells & Singleton-cell comparison \\
Firm-buyer-item triples & \BPutgTripleNUTG{} (\BPutgTripleCountUTG{} triples) & Winning bids & Within-supplier pricing test \\
\bottomrule
\end{tabular}
\begin{tablenotes}[flushleft]\footnotesize
\item \textit{Notes:} POI denotes purchase-offer-item; item-month cells are item-by-year-month cells. Price regressions use accepted winning bids. Classifier validation is conducted at the purchase-order/tender-notice level (Online Appendix~A). The firm-buyer-item triple sample is a strict subset of urgent winning bids and identifies within-supplier pricing, not supplier reallocation.
\end{tablenotes}
\end{threeparttable}
\end{table}
